% Seção: Implementação
% Dono: Thiago e Arthur · Orçamento: 1 página(s) · docs/fontes/plano-de-projeto.md §14.1
% Conteúdo: Núcleo de grafo e heap, organização do código e interface
% Condensada em 20/11 para o orçamento de 15 páginas; o diagrama de sequência da página
% Início (figuras/sequencia-p0-ui.png) segue em docs/diagramas/.
\section{Implementação}\label{sec:impl}

O núcleo (\texttt{src/transporte/grafo/}) não usa biblioteca de grafo nem o \texttt{heapq}. O
grafo guarda os atributos de vértices e arestas em listas paralelas indexadas por identificadores
estáveis, e a lista de adjacência de cada vértice contém os identificadores das arestas que saem
dele, o que permite arestas paralelas (viagem e caminhada entre as mesmas paradas) e devolve, no
caminho, o tipo de cada passo. Vértices e arestas são desativados por uma marcação em $O(1)$ e
as buscas os ignoram, o que permite simular falhas sem reconstruir o grafo. A fila de prioridade
é um heap binário de mínimo próprio. No feed real, um Dijkstra completo leva 250~ms em $G_P$ e
777~ms em $G_E$ do pico da manhã, com distâncias iguais às do NetworkX em 20 origens; o heap faz
100 mil inserções e extrações em 1,5~s (Intel Core 5 120U, Python 3.12). O grafo só aceita pesos
não negativos e nunca remove vértices ou arestas.

O gancho \texttt{.githooks/pre-commit} analisa a árvore sintática dos módulos de produção para
detectar imports proibidos e executa o Ruff. A interface é uma aplicação Streamlit com uma
página por pergunta: a página Início carrega a pasta ou ZIP do GTFS, reutiliza o resultado do
leitor pelo cache e copia a rede para cada sessão; as demais recebem a rede e desenham os
resultados em camadas Folium. No serviço de 04/11 (macOS arm64, Python 3.12.13, três
repetições), a carga levou $7{,}827 \pm 0{,}034$~s com cache vazio e $0{,}671 \pm 0{,}017$~s com
cache preenchido, e montar o resumo de $G_P$ e $G_E$, $0{,}440 \pm 0{,}001$~s.
